\lesson{2}{Mar 18 2022 Fri (17:40:16)}{Completing the Square}{Unit 4}
\label{les_2_unit_4:completing_the_square}

\begin{definition}[Completing the Square]
  \label{def:completing_the_square}

  Completing the square means writing a quadratic in the form of a squared
  bracket and adding a constant (if necessary).
\end{definition}

\begin{example}
  \label{exm:completing_the_square}

  Make $x^{2} - 4x - 13 = 0$ into a perfect square.

  \begin{equation}
    \label{spt:completing_the_square}
    \begin{split}
      x^{2} - 4x - 13 &= 0 \\
      x^{2} - 4x = 13 \\
    \end{split}
  .\end{equation}

  Now, we're going to find a value ($c$), that will create a perfect square
  trinomial. To find our new value, we're going to take:

  \begin{equation}
    \label{spt:finding_our_b}
    \begin{split}
      \left(-\frac{4}{2}\right)^{2} = 4
    \end{split}
  .\end{equation}

  Which gives us:

  \begin{equation}
    \label{spt:completing_the_square_2}
    \begin{split}
      (x^{2} - 4x + 4) = 13 + 4
    \end{split}
  .\end{equation}

  Which can be re-written as a perfect square:

  \begin{equation}
    \label{spt:completing_the_square_3}
    \begin{split}
      (x - 2)^{2} = 17
    \end{split}
  .\end{equation}
\end{example}

\newpage
